\documentclass[aspectratio=169, onlytextwidth,  notheorems, sOuRcEs, fleqn, leqno, ]{RUCPresentation}

%\usepackage{pgfforeach}
% \usepackage{foreach}

\newcounter{chapter}


\usepackage{tikz}
\usetikzlibrary{calc}
\usetikzlibrary{angles}
\usetikzlibrary{positioning}

\usetikzlibrary{arrows.meta}
\usetikzlibrary{shapes.geometric}
\usetikzlibrary{overlay-beamer-styles}

\ProvideDocumentCommand\pgfmathsetnewmacro{ m m }{ \newcommand*{#1}{} \pgfmathsetmacro{#1}{#2} }
\ProvideDocumentCommand\pgfmathsetnewlengthmacro{ m m }{ \newcommand*{#1}{} \pgfmathsetlengthmacro{#1}{#2} }

\tikzset{>=Stealth} % for good looking LaTeX arrow head
% \tikzset{overlay, remember picture}



\usepackage{myMathQA}

\usepackage{ulem}

\usepackage{graphicx}
\usepackage{siunitx}
\usepackage{subcaption}

\usepackage{hyperref}

% \usepackage[style=authoryear]{biblatex}
\usepackage{biblatex}
\addbibresource{mscbib.bib}

\DeclareSourcemap{
    \maps[datatype=bibtex]{
        \map[overwrite]{
            \step[fieldset=note, null]
        }
        \map[overwrite]{
            \step[fieldset=ppn_gvk, null]
        }
        \map[overwrite]{
            \step[fieldset=file, null]
        }
        \map[overwrite]{
            \step[fieldset=doi, null]
        }
        \map[overwrite]{
            \step[fieldset=issn, null]
        }
        \map[overwrite]{
            \step[fieldset=isbn, null]
        }
        \map[overwrite]{
            \step[fieldset=pagetotal, null]
        }
    }
}

% \DeclareCiteCommand{\beamerfootcite}
%   {}
%   {\footnote[frame]{\printbibliography[keys=\thefield{entrykey}]}}
%   {\multicitedelim}
%   {\usebibmacro{postnote}}

\ProvideDocumentCommand{\beamerfootcite}{ m }{\footnote[frame]{\ \fullcite{#1}}}
\renewcommand*{\nameyeardelim}{\addcomma\space} % cleaner authoryear style

% \AtBeginSection{ \frame{\sectionpage} }
    % \begin{frame}{\insertsection}
    %     \sectionpage
    % \end{frame}

% \usepackage{animate}



\def\k{\raisebox{1ex}{\rotatebox{180}{k}}}
% \newsavebox\tempboxbeamer

\usetheme{RUCTHEME}
% \usetheme{Warsaw}

\makeatletter
\renewcommand\@makefnmark{\hbox{\@textsuperscript{\usebeamercolor[fg]{footnote mark}\usebeamerfont*{footnote mark}[\@thefnmark]}}}
% \renewcommand\@makefntext[1]{{\usebeamercolor[fg]{footnote mark}\usebeamerfont*{footnote mark}[\@thefnmark]}\usebeamerfont*{footnote} #1}
\makeatother

\hypersetup{colorlinks = false, citecolor=TextColor}
\AtBeginBibliography{\small\color{TextColor}}


\def\dumbmatrixtwobytwo[#1,#2]{%
    \count6=#1\relax \count7=#2\relax%
    \count2=#1\relax \count3=#2\relax \count4=#1\relax \count5=#2\relax%
    \multiply\count2 by \count6\relax \multiply\count4 by \count7\relax%
    \multiply\count3 by \count6\relax \multiply\count5 by \count7\relax%
    \ensuremath{%
        \left[%
            \begin{matrix}%
                \number\count2 & \number\count4 \\%
                \number\count3 & \number\count5%
            \end{matrix}%
        \right]%
    }
}


\def\dumbmatrixthreebythree[#1,#2,#3]{%
    \count10=#1\relax \count11=#2\relax \count12=#3\relax
    \count1=#1\relax \count4=#1\relax \count7=#1\relax%
    \count2=#2\relax \count5=#2\relax \count8=#2\relax%
    \count3=#3\relax \count6=#3\relax \count9=#3\relax%
    \multiply\count1\count10\relax \multiply\count4\count11\relax \multiply\count7\count12\relax%%
    \multiply\count2\count10\relax \multiply\count5\count11\relax \multiply\count8\count12\relax%%
    \multiply\count3\count10\relax \multiply\count6\count11\relax \multiply\count9\count12\relax%%
    \ensuremath{%
        \left[%
            \begin{matrix}%
                \number\count1 & \number\count4 & \number\count7 \\%
                \number\count2 & \number\count5 & \number\count8 \\%
                \number\count3 & \number\count6 & \number\count9%
            \end{matrix}%
        \right]%
    }
}

\usepackage{multirow}
\newcounter{reaction}
\ProvideDocumentCommand\reactionterm{ o m m m m }{% --> Function for reaction lines
    \stepcounter{reaction} %
    \IfNoValueTF{#1}
    { \arabic{reaction} & #3 & #4 & \(t^+_{\arabic{reaction}}\of{\vec{x}}\) & #2 & \(#5\) \\ \hline }
    { \multirow{2}*{\arabic{reaction}} & #3 & #4 & \(t^+_{\arabic{reaction}}\of{\vec{x}}\) & #2 &  \multirow{2}*{\(#5\)} \\
                                       & #4 & #3 & \(t^-_{\arabic{reaction}}\of{\vec{x}}\) & #1 &                             \\ \hline }
}

\ProvideDocumentEnvironment{reactiontable}{}%
{ \begin{center} \begin{tabular}{c|r@{\ \(\to\)\ }lr@{\ \(=\)\ }lc} %
    \(I\) & \multicolumn{2}{c}{`Reaction'} %
          & \multicolumn{2}{c}{`Rate Law'} %
          & \(\vec{r}_I^{\mathrm{T}}\) \\\hline }%
{ \end{tabular} \end{center} \setcounter{reaction}{0} }


% Defining a new coordinate system for the page:
%
% --------------------------
% |(-1,1)    (0,1)    (1,1)|
% |                        |
% |(-1,0)    (0,0)    (1,0)|
% |                        |
% |(-1,-1)   (0,-1)  (1,-1)|
% --------------------------


% For beamer title
\newcommand\makebeamertitle{\frame{\maketitle}}%

% For sections?
\graphicspath{{./images/}}%{../Thesis/images/}{../Thesis/images/phase/}{../Thesis/images/ode/}{../Thesis/images/fpe/}}


\def\cfp#1{c_#1^*}


\title{%
    Some Stuff Related to my Thesis
}
\subtitle{Presentation For The MSc Thesis Exam}

% \author{by Anakin}
\author[Anakin]{Anakin Diermayr\inst{1}}% \newline \texttt{anakin@ruc.dk} }
\institute[RUC]{
    \inst{1} Roskilde University
}

%\mode<presentation>


\begin{document}


\makebeamertitle


% \section*{Outline}
% \begin{frame}[allowdisplaybreaks=4]{Outline}
%
%     \tableofcontents[pausesections]
%
% \end{frame}

% \begin{frame}{Workshopping}
%
%     \begin{enumerate}
%         \item Welcome/Introduction
%         \item Brief reminder of the themes/goals
%         \item Missing Details
%             \begin{enumerate}
%                 \item Martingales?
%                 \item Birth-Death Processes < -- --
%                 \item How it applies with a 1 dimensional version of the model
%                     I used.
%             \end{enumerate}
%     \end{enumerate}
%
% \end{frame}


\section{Introduction}

\begin{frame}{Introduction}

    % \vspace{2em}


    \begin{columns}[T]

        \column{0.50\framewidth}
        \begin{itemize}
            \item Welcome to my Presentation
            \item My name is Anakin\uncover<2->{... yes, like the StarWars}
            \item<3-> BSc  at RUC, combination Molecular Biology, Mathematics
        \end{itemize}

        \column<4->{0.50\framewidth}
        My Original Ambition:
        \begin{itemize}
            % \item Find a way to formalizing Waddington's `epigentic landscape'.%\beamerfootcite{Waddington2014}
            \item Formalizing Waddington's `epigentic landscape'.%\beamerfootcite{Waddington2014}
            \item An excuse to `balls deep' into stochasticity
        \end{itemize}
    \end{columns}

    \begin{alertblock}<5->{Alert}
        This presentation will not directly cover the material in my Thesis.\par
        We will cover motivation/inspiration and tie it to results.
    \end{alertblock}

\end{frame}

% \item I believe there is an inherent relationship between
%     genetic phenotypes and enviromental pressures that
%     only statistical manifolds can represent


% TODO:
\subsection{ Motivation }
\begin{frame}{ Motivation }
    % \vspace{1em}

    \begin{columns}

        \column{0.45\framewidth}
        \begin{itemize}[<+->]
            \item The Physical Analogy
                \begin{itemize}[<.->]
                    \item Potential Energy Landscapes
                    \item Valleys as Stable phenotypes
                \end{itemize}
            % \item I considered Information Geometry as a viable route
            \item Information Geometry Route
                \begin{itemize}[<.->]
                    \item Statistical Geometry %would introduce a natural
                    \item Frobenius Algebras/Manifolds% have Appropriate Requirements
                    \item Natural construction from random variables
                \end{itemize}
            \item However, I lacked a lot of the pre-requisites educationally
        \end{itemize}


        \column{0.60\framewidth}

        \begin{figure}[htbp]
        \centering
        \includegraphics[width=0.7\textwidth]{epigenetic.png}
        \caption{Part of an epigentic landscape\beamerfootcite{Waddington2014}.}
        \end{figure}



        % \includegraphics[width=0.5\framewidth]{placehold.pdf}

    \end{columns}


\end{frame}

\subsection{The Research Question}
\begin{frame}[t]{The Research Question}

    % \vspace{2em}

        % Why I chose this
        % \begin{enumerate}[<+->]
        %     \item It's a tangible problem
        %     \item I wanted to find an excuse to go `balls deep' into stochasticity
        %     \item I believe there is an inherent relationship between
        %         genetic phenotypes and enviromental pressures that
        %         only statistical manifolds can represent
        % \end{enumerate}

    % \item I needed something more concrete to begin working with.
    \begin{itemize}[<+->]
        \item After \alt<2->{\xout{some} \emph{a lot} of}{some} deliberation,
            % and an expectation of data from a molecular biologist,
            %I settled on
            \begin{block}<+->{The Reseach Question:}
                ``How can stochastic gene copy number dynamics underlying bacterial
                heteroresistance be accurately captured using a Fokker-Planck formulation,
                and what insights does this stochastic framework provide beyond deterministic ODE models.''
            \end{block}
    \end{itemize}


    \tikz[overlay, remember picture] \node<2-> at (frame cs:0.35,-0.4) {%
        \begin{tikzpicture}
            \pgfmathsetnewmacro\genelength{35}
            \pgfmathsetnewmacro\geneheight{2}

            \pgfmathsetnewmacro\translength{10}
            \pgfmathsetnewmacro\transheight{2}

            \pgfmathsetnewmacro\reactionthickness{1}
            \pgfmathsetnewmacro\dnawidth{4}

            \pgfmathsetnewmacro\shortendist{6}
            \pgfmathsetnewmacro\bendstrength{1.05} % √(1²+(0.2-0.75)²)

            \pgfmathsetnewmacro\geneshift{5}

            \providecommand*\genecolor{red!80!black}
            \providecommand*\transcolor{green!80!black}

            \useasboundingbox (-4,-1.25) rectangle (5.5,1.35);
            \coordinate (O) at (0,0);

            \node[circle] (west) at (-3.5, 0) {};
            \node[circle] (east) at ( 4.5, 0) {};

            \node[anchor=south east] at (west) {\(5'\)};
            \node[anchor=south west] at (east) {\(3'\)};

            \draw[line width = \dnawidth] (west.center) to (east.center);

            \draw[line width = \dnawidth] (east.north) to (east.south);
            \draw[line width = \dnawidth] (west.north) to (west.south);


            \coordinate (g1) at (-2,0); \coordinate (g2) at ( 0,0);
            \coordinate (g3) at ( 3,0);

            \coordinate (t1) at (-3,0); \coordinate (t2) at (-1,0);
            \coordinate (t3) at ( 1,0);
            \coordinate (t4) at ( 2,0); \coordinate (t5) at ( 4,0);

            \draw[line width = \dnawidth+1, BackgroundColor, dashed] (t3) -- (t4);
            \draw[line width = \dnawidth+1, BackgroundColor] ([xshift=0.37cm]t3) to ([xshift=-0.43cm]t4);

            \foreach \POINT in {1,2,3}{
                \node[rectangle, fill=\genecolor, draw=\genecolor, minimum width=\genelength, minimum height=\geneheight] at (g\POINT) {};
            }

            \foreach \POINT in {1,2,3,4,5}{
                \node[rectangle, fill=\transcolor, draw=\transcolor, minimum width=\translength, minimum height=\transheight] at (t\POINT) {};
            }

            \node[above, yshift=5] (x1) at (g1) {\(X\)};
            \node[above, yshift=5] (x2) at (g2) {\(X\)};
            \node[above, yshift=5] (x3) at (g3) {\(X\)};

            \node  (B) at ( 4.75, 0.85) {\(B\)};
            \node (B2) at ( 0.50, 0.85) {\(B\)};

            \node  (N) at ( 2.00,-0.75) {\(N\)};
            \node (N2) at (-2.00, 1.00) {\(N\)};
            \node (N3) at ( 2.75, 1.00) {\(N\)};

            \node  (y) at ( 4.00,-0.75) {\(Y\)};
            \node (y2) at ( 1.50, 0.30) {\(Y\)};
            \node (y3) at ( 0.75,-0.75) {\(Y\)};

            \begin{scope}[line width=\reactionthickness, shorten <=\shortendist, shorten >=\shortendist]
                \draw[->] ([yshift=\geneshift]g1)
                    to[bend left] node[midway, sloped] (k1path) {}
                    ([yshift=\geneshift]g2);
                \draw[->] ([yshift=-\geneshift]g2)
                    to[bend left] node[midway] (k-1path) {}
                    ([yshift=-\geneshift]g1);
            \end{scope}

            \node[above] at (k1path) {\(k_1\)};
            \node[below] at (k-1path) {\(k_{-1}\)};

            \draw[line width=\reactionthickness, shorten >=6pt]
                (N2) to[bend right] ([yshift=2.3]k1path.center);

            % \draw[->,line width=\reactionthickness, shorten <=5.5pt]
            %     ([yshift=-2.6]k-1path.center) to[bend right] (N3);



            \begin{scope}[line width = \reactionthickness]
                \draw[shorten >=5pt, shorten <=4pt ]
                    (B.center) -- (x3.center)
                    node[pos=0.5, sloped] (midbx) {\(\bigtimes\)};
                \draw[shorten >=5pt, shorten <=4pt ]
                    (B2.center) -- (y2.center)
                    node[pos=0.5, sloped] (midby) {\(\bigtimes\)};

                \draw[->] (midbx.center) to[in=0, out=105, looseness=1.5] (N3);
                \draw[->] (midbx.center) to[in=160, out=105, looseness=3] (B);

                \draw[->] (midby.center) to[in=180, out=71, looseness=1.5] (N3);

                % \draw[->] (midxb.center) -- ++(105:0.5);
                % \draw[->] (midxb.center) to[bend right] (N2.east);

                \draw[->] (N) .. controls ++( 45:\bendstrength) and ++(135:\bendstrength) .. (y)
                    node[midway, below] {\(k_2\)};
                \draw[->] (y3) -- (N) node[midway, above] {\(k_5\)};
            \end{scope}

            \node[anchor=north west] at (midbx) {\(k_3\)};
            \node[anchor=west] at ([yshift=2,xshift=1]midby) {\(k_4\)};


        \end{tikzpicture}
    };

    \tikz[overlay, remember picture] \node<2-> at (frame cs:-0.65,-0.35) {
        \begin{tikzpicture}[ scale = 1,
            draw = black!80,
            thick,
            shorten >=2pt,
            overlay
            ]

            \node[draw=black, circle] (c1) at (-1, 0) {\(c_1\)};
            \node[draw=black, circle] (c2) at ( 1, 0) {\(c_2\)};

            \node[draw=black] (m) at ( -1, -1.20) {\(m\)};
            \node[] (source) at ( -2, -1.20) {\(m_0\)};

            \draw[->] (c2) [bend left]to node[midway, below] {\(w_2\)} (c1);
            \draw[->] (c1) [bend left]to node[midway, above] {\(w_1\)} (c2);

            \draw[->] (c1)  [loop left]to node[midway, left]  {\(k_1\)} (c1);
            \draw[->] (c2) [loop right]to node[midway, right] {\(k_2\)} (c2);

            \draw[-|] (m) to node[pos=0.35] (mm) {} (c1);
            \draw[-|, rounded corners] (c2) |- node[midway,anchor=south west] {\(q_m\)} (m);
            % \draw[-|] (c2) to[out=270, in=0, looseness=2] node[midway,anchor=north west] {\(q_m\)} (m);

            \node[anchor=east] at (mm) {\(q_1\)};

            \draw[->] (source) -- (m);

    \end{tikzpicture}};

    % \begin{tikzpicture}[overlay, remember picture]
    %
    %     \node<2-> at (frame cs: 0, -0.40) {%
    %             % \includegraphics[width=0.65\framewidth]{cgn.png}
    %         };
    %
    % \end{tikzpicture}

\end{frame}

% \section{ Why this stuff goes hard }
\section{ Real World Examples }


% \subsection{}

% \begin{frame}{ Real World Examples }
%
%
%     \begin{tikzpicture}[overlay, remember picture]
%         \node<1>[fill] (a) at (frame cs: 0, 0.18) {%
%                 \includegraphics[width=0.85\framewidth]{phenotypespace.png}
%             };
%
%
%
%         \node at (frame cs: 0, -0.5) {%
%                 Phase portraits characterize the genotype–phenotype map\beamerfootcite{Jaeger2014}
%             };
%     \end{tikzpicture}
%
%     % \begin{figure}[htbp]
%     % \centering
%     % \includegraphics[width=0.7\textwidth]{phenotypespace.png}
%     % \caption{Phase portraits characterize the genotype–phenotype map\beamerfootcite{Jaeger2014}}
%     % \end{figure}
%
% \end{frame}

\begin{frame}{ Dongya Jia et al. }


    \begin{tikzpicture}[overlay, remember picture]

        \begin{scope}[fill]
            \node[fill] (a) at (frame cs: -0.52, 0.2) {
                    \includegraphics[width=0.48\textwidth]{cancers-09-00070-g002--1.png}
                };
            \node[fill] (b) at (frame cs: 0.35, 0.2) {
                    \includegraphics[width=0.48\textwidth]{cancers-09-00070-g002--2.png}
                };
            \node[fill] (c) at (frame cs: 0.85, 0.2) {
                    \includegraphics[height=0.35\framewidth]{cancers-09-00070-g002--3.png}
                };
        \end{scope}

        \node[text width=0.8\framewidth] at (frame cs: -0.05, -.38) {
                Schematic illustration of the quasi-potential landscape for epithelial-to-mesenchymal transition (EMT) in 3-dimensional space\beamerfootcite{Jia2017}
            };

    \end{tikzpicture}


\end{frame}


\begin{frame}{ Michael Levin }

    \begin{tikzpicture}[overlay, remember picture]
        \node<1>[fill] (a) at (frame cs: 0, 0.18) {%
                \includegraphics[width=0.85\framewidth]{cellregeneration.png}
            };

        \node[text width = 0.9\textwidth] at (frame cs: 0, -0.46) {%
                Kidney tubules form with roughly correct diameter in newts despite polyploidy, which changes cell size. \beamerfootcite{Levin2021}
            };
    \end{tikzpicture}

\end{frame}

% \begin{frame}{ Levin }
%
%     % \begin{tikzpicture}[overlay, remember picture]
%     %     \node<1>[fill] (a) at (frame cs: 0, 0.18) {%
%     %             \includegraphics[width=0.85\framewidth]{cellregeneration.png}
%     %         };
%     %
%     %     \node[text width = 0.8\textwidth] at (frame cs: 0, -0.5) {%
%     %             Kidney tubules form with roughly correct diameter in newts despite polyploidy, which changes cell size. \beamerfootcite{Levin2021}
%     %         };
%     % \end{tikzpicture}
%
% \end{frame}



\section{ Results }


\begin{frame}{ A Question Emerges }


    % \vspace{1em}

    \begin{onlyenv}<1>

    \begin{itemize}
        \item How does any of this relate to the results?
    \end{itemize}

    \end{onlyenv}

    \begin{onlyenv}<2->
        \begin{proposition}[name={Stochastic Stability Decay}] \label{prp.ssd}

            %which we can represent with \gls{fpe}
            Assume \(M\) is a manifold denoting a state space with an unstable origin,
            whose {evolution rule} is \(\varphi:\R\times M\to M\),
            \(t\in\R\), and that \(N\subset M\) is a submanifold that includes the origin.

            % with an attracting point \(\vec{x}^*\in N\backslash\{\vec{0}\}\).
            Additionally assume we have system of SDEs,
            which involves an isolated autocatalytic process,
            \begin{equation*}
                \hfill
                \delta \vec{x}\of t
                = \vec{A}\br{\vec{x}\of t, t}
                + \sqrt{\bfm B\br{\vec{x}\of t, t}} \ \delta{\vec{W}\of{t\st \vec{x}}}
                \hfill
            \end{equation*}
            Such that the drift term \(\vec{A}:N\to\R^n\) is the vector field
            associated with the flow \(\varphi_t\of N\).
            % homeomorphic to \(N\).

            If \(\varphi_t\of{N}\) can be partitioned into a trapping region,
            \(N\backslash\{\vec{0}\}\), and the unstable point \(\{\vec{0}\}\).
            Then given a large enough time interval,
            the stochastic system will inevitably decay from stability.

        \end{proposition}
    \end{onlyenv}

\end{frame}


\subsubsection{ODE Model Potentials}
\begin{frame}{ Stochastic Processes }

    \begin{tikzpicture}[overlay, remember picture]
        \node<1> at (frame cs: -0.45, 0.08) {%
                \includegraphics[width=0.55\framewidth, page=1]{ssa:ODE8Par3Var:I34:C1T17476639695953887:2dHist.pdf}
            };
        \node<1> at (frame cs: 0.45, 0.08) {%
                \includegraphics[width=0.50\framewidth, page=2]{ssa:ODE8Par3Var:I34:C1T17476639695953887:2dHist.pdf}
            };

        \node[text width = 0.9\textwidth] at (frame cs: 0, -0.66) {%
                The stochastic process logically equivalent to the ODE model
                outlined in Section 6.1
            };
    \end{tikzpicture}

\end{frame}

\begin{frame}{ Analytical/Numeric Processes }

    \begin{onlyenv}<1>

    \begin{columns}
        \column{0.5\framewidth}
                \includegraphics[width=1\textwidth, page=1]{potential_ODE8Par3Var--1.pdf}
                The distribution of the system when \(m_0 = 0\)


        \column{0.5\framewidth}
                \includegraphics[width=1\textwidth, page=1]{potential_ODE8Par3Var--100.pdf}
                The distribution of the system when \(m_0 = 10\)

    \end{columns}

    \end{onlyenv}

    \begin{onlyenv}<2>

        \begin{columns}
            \column{0.5\framewidth}
            \includegraphics[width=1\textwidth, page=1]{potential_ODE8Par3Var--Potential--M0=0--1.pdf}
                The potential landscape of the system when \(m_0 = 0\)


            \column{0.5\framewidth}
            \includegraphics[width=1\textwidth, page=1]{potential_ODE8Par3Var--Potential--M0=10--1.pdf}
            % \includegraphics[width=1\textwidth, page=1]{potential_ODE8Par3Var--Potential--100.pdf}
                The potential landscape of the system when \(m_0 = 10\)

        \end{columns}

    \end{onlyenv}

\end{frame}


\subsubsection{My Model Potentials}
\begin{frame}{ Stochastic Processes }

    \begin{tikzpicture}[overlay, remember picture]
        \node<1> at (frame cs: -0.45, 0.08) {%
                \includegraphics[width=0.55\framewidth, page=1]{ssa:Novel5Par3Var:I26:C1T17475033100262814:2dHist.pdf}
            };
        \node<1> at (frame cs: 0.45, 0.08) {%
                \includegraphics[width=0.50\framewidth, page=2]{ssa:Novel5Par3Var:I26:C1T17475033100262814:2dHist.pdf}
            };

        \node[text width = 0.9\textwidth] at (frame cs: 0, -0.66) {%
                The stochastic process of my own model outlined in Section 6.2
            };
    \end{tikzpicture}

\end{frame}

\begin{frame}{ Analytical/Numeric Processes }

    \begin{onlyenv}<1>

        \begin{columns}
            \column{0.5\framewidth}
            \includegraphics[width=1\textwidth, page=1]{potential_Novel5Par3Var--B=0N--1.pdf}
                The distribution of the system when \(B = 0\)


            \column{0.5\framewidth}
            \includegraphics[width=1\textwidth, page=1]{potential_Novel5Par3Var--B=N--1.pdf}
                The distribution of the system when \(B = N\)

        \end{columns}

    \end{onlyenv}

    \begin{onlyenv}<2>

    \begin{columns}
        \column{0.5\framewidth}
                \includegraphics[width=1\textwidth, page=1]{potential_Novel5Par3Var--Potential--B=0N--1.pdf}
                The potential landscape of the system when \(B = 0\)


        \column{0.5\framewidth}
                \includegraphics[width=1\textwidth, page=1]{potential_Novel5Par3Var--Potential--B=N--1.pdf}
                The potential landscape of the system when \(B = N\)

    \end{columns}

\end{onlyenv}

\end{frame}






% \subsection{What I did}
% \begin{frame}{The process}
%
%     \begin{itemize}
%         \item I had to cover a lot of novel ground
%         \item
%     \end{itemize}
%
% \end{frame}
% \section{ Meat 'n Potatoes }
% \subsection{Background Elaboration}

% \begin{frame}{ Birth Death Processes }
%
%     \begin{columns}[T]
%
%         \column{0.5\framewidth}
%
%         \begin{itemize}
%             \item Discrete Stochastic Processes
%                 \begin{itemize}
%                     \item Jump Processes
%                     \item Birth-Death Processes
%                 \end{itemize}
%             \item Continuous Stochastic Processes
%         \end{itemize}
%
%         \column{0.5\framewidth}
%
%         \begin{itemize}
%             \item
%         \end{itemize}
%
%     \end{columns}
%
%
% \end{frame}

% \subsection{Gene stuff}

% \begin{frame}{ Recapping the Model }
%
%
%     \vspace{1em}
%     \begin{columns}
%
%         \column{0.4\framewidth}
%
%         \begin{onlyenv}<1>
%
%             \begin{itemize}
%                 \item[\(X\)] represents the gene copy number
%                 \item[\(Y\)] represents the protective resistance proteins
%                 \item[\(N\)] represents the available nutrients in the system
%                 \item[\(B\)] represents the antibiotic added to the system
%                 \item[\(\emptyset\)] represents elements outside the scope of the system
%             \end{itemize}
%
%         \end{onlyenv}
%
%         \begin{onlyenv}<2->
%             The reactions:
%             \begin{align}
%                 X + N &\xrightleftharpoons[k_{-1}]{k_1} 2 X \\
%                 X + N &\xrightarrow{k_2} X + Y \\
%                 X + B &\xrightarrow{k_3} B \\
%                 Y + B &\xrightarrow{k_4} N \\
%                 Y     &\xrightarrow{k_5} N \\
%                 % \uncover<3->{\emptyset &\xrightarrow{k_b} B \\ }
%                 \notag
%             \end{align}
%         \end{onlyenv}
%
%         \column{0.50\framewidth}
%
%     \end{columns}
%
%     \tikz[overlay, remember picture] \node at (frame cs:0.35,0) {%
%         \begin{tikzpicture}
%             \pgfmathsetnewmacro\genelength{35}
%             \pgfmathsetnewmacro\geneheight{2}
%
%             \pgfmathsetnewmacro\translength{10}
%             \pgfmathsetnewmacro\transheight{2}
%
%             \pgfmathsetnewmacro\reactionthickness{1}
%             \pgfmathsetnewmacro\dnawidth{4}
%
%             \pgfmathsetnewmacro\shortendist{6}
%             \pgfmathsetnewmacro\bendstrength{1.05} % √(1²+(0.2-0.75)²)
%
%             \pgfmathsetnewmacro\geneshift{5}
%
%             \providecommand*\genecolor{red!80!black}
%             \providecommand*\transcolor{green!80!black}
%
%             \useasboundingbox (-4,-1.25) rectangle (5.5,1.35);
%             \coordinate (O) at (0,0);
%
%             \node[circle] (west) at (-3.5, 0) {};
%             \node[circle] (east) at ( 4.5, 0) {};
%
%             \node[anchor=south east] at (west) {\(5'\)};
%             \node[anchor=south west] at (east) {\(3'\)};
%
%             \draw[line width = \dnawidth] (west.center) to (east.center);
%
%             \draw[line width = \dnawidth] (east.north) to (east.south);
%             \draw[line width = \dnawidth] (west.north) to (west.south);
%
%
%             \coordinate (g1) at (-2,0); \coordinate (g2) at ( 0,0);
%             \coordinate (g3) at ( 3,0);
%
%             \coordinate (t1) at (-3,0); \coordinate (t2) at (-1,0);
%             \coordinate (t3) at ( 1,0);
%             \coordinate (t4) at ( 2,0); \coordinate (t5) at ( 4,0);
%
%             \draw[line width = \dnawidth+1, BackgroundColor, dashed] (t3) -- (t4);
%             \draw[line width = \dnawidth+1, BackgroundColor] ([xshift=0.37cm]t3) to ([xshift=-0.43cm]t4);
%
%             \foreach \POINT in {1,2,3}{
%                 \node[rectangle, fill=\genecolor, draw=\genecolor, minimum width=\genelength, minimum height=\geneheight] at (g\POINT) {};
%             }
%
%             \foreach \POINT in {1,2,3,4,5}{
%                 \node[rectangle, fill=\transcolor, draw=\transcolor, minimum width=\translength, minimum height=\transheight] at (t\POINT) {};
%             }
%
%             \node[above, yshift=5] (x1) at (g1) {\(X\)};
%             \node[above, yshift=5] (x2) at (g2) {\(X\)};
%             \node[above, yshift=5] (x3) at (g3) {\(X\)};
%
%             \node  (B) at ( 4.75, 0.85) {\(B\)};
%             \node (B2) at ( 0.50, 0.85) {\(B\)};
%
%             \node  (N) at ( 2.00,-0.75) {\(N\)};
%             \node (N2) at (-2.00, 1.00) {\(N\)};
%             \node (N3) at ( 2.75, 1.00) {\(N\)};
%
%             \node  (y) at ( 4.00,-0.75) {\(Y\)};
%             \node (y2) at ( 1.50, 0.30) {\(Y\)};
%             \node (y3) at ( 0.75,-0.75) {\(Y\)};
%
%             \begin{scope}[line width=\reactionthickness, shorten <=\shortendist, shorten >=\shortendist]
%                 \draw[->] ([yshift=\geneshift]g1)
%                     to[bend left] node[midway, sloped] (k1path) {}
%                     ([yshift=\geneshift]g2);
%                 \draw[->] ([yshift=-\geneshift]g2)
%                     to[bend left] node[midway] (k-1path) {}
%                     ([yshift=-\geneshift]g1);
%             \end{scope}
%
%             \node[above] at (k1path) {\(k_1\)};
%             \node[below] at (k-1path) {\(k_{-1}\)};
%
%             \draw[line width=\reactionthickness, shorten >=6pt]
%                 (N2) to[bend right] ([yshift=2.3]k1path.center);
%
%             % \draw[->,line width=\reactionthickness, shorten <=5.5pt]
%             %     ([yshift=-2.6]k-1path.center) to[bend right] (N3);
%
%
%
%             \begin{scope}[line width = \reactionthickness]
%                 \draw[shorten >=5pt, shorten <=4pt ]
%                     (B.center) -- (x3.center)
%                     node[pos=0.5, sloped] (midbx) {\(\bigtimes\)};
%                 \draw[shorten >=5pt, shorten <=4pt ]
%                     (B2.center) -- (y2.center)
%                     node[pos=0.5, sloped] (midby) {\(\bigtimes\)};
%
%                 \draw[->] (midbx.center) to[in=0, out=105, looseness=1.5] (N3);
%                 \draw[->] (midbx.center) to[in=160, out=105, looseness=3] (B);
%
%                 \draw[->] (midby.center) to[in=180, out=71, looseness=1.5] (N3);
%
%                 % \draw[->] (midxb.center) -- ++(105:0.5);
%                 % \draw[->] (midxb.center) to[bend right] (N2.east);
%
%                 \draw[->] (N) .. controls ++( 45:\bendstrength) and ++(135:\bendstrength) .. (y)
%                     node[midway, below] {\(k_2\)};
%                 \draw[->] (y3) -- (N) node[midway, above] {\(k_5\)};
%             \end{scope}
%
%             \node[anchor=north west] at (midbx) {\(k_3\)};
%             \node[anchor=west] at ([yshift=2,xshift=1]midby) {\(k_4\)};
%
%
%         \end{tikzpicture}
%     };
%
%
% \end{frame}





% TODO:

% \begin{frame}{The Corollary}
%
%     \begin{corollary} \label{con:timetodestable}
%
%         Assuming that we have a system that fufills the properties of a stochastic
%         stability decay.
%         And assuming that \(Z\) is a standard normal random variable.
%         Then, the median interval of time, \(t-s\), \(s<t\in[0,\infty[\),
%         needed to hit the origin, will be dependent on the distance
%         \(\mu\of{0, x^*}= L\) and the diffusion term, such that
%         \begin{equation*}
%             t-s = \Delta t = \pa{
%                 {L \over z_\alpha \sqrt{|\nabla \cdot \bfm B\of{\vec{x}^*}|}}
%             }^2
%         \end{equation*}
%         will give a \(100\pa{1-\alpha}\) confidence interval for the
%         system to decay to zero.
%
%     \end{corollary}
%
% \end{frame}


% \begin{frame}{Does this even make sense as a claim?}
%
%     \begin{align*}
%         \nabla\cdot \mat{B} &= \nabla \cdot \bmatform[ -k_{-1} x + k_1 y,  k_{-1} x - k_1 y; k_{-1} x - k_1 y, -k_{-1} x + k_1 y ]
%          = \bmatform[ 0, 0]
%     \end{align*}
%
%     \begin{align*}
%         \nabla\cdot \mat{B} &= \nabla \cdot \bmatform[
%         \pa{n +  c_1 + c_2} k_1 c_1 +  w_1 c_1 + w_2 c_2 ,
%         - \pa{ w_1 c_1 + w_2 c_2  };
%         - \pa{ w_1 c_1 + w_2 c_2  },
%         \pa{n +  c_1  + c_2} k_2 c_2 + w_1 c_1 + w_2 c_2
%         ] \\
%         &= \bmatform[
%         \pa{n_1 + 2c_1 +  c_2}k_1,
%         \pa{n_2 +  c_1 + 2c_2}k_2
%         ] \\
%        \implies |\nabla\cdot \mat{B}\of{\vec{x}^*}| &= \sqrt{
%            (\pa{n_1 + 2c_1 +  c_2}k_1)^2 + (\pa{n_2 +  c_1 + 2c_2}k_2)^2
%        }
%     \end{align*}
%
% \end{frame}

% % TODO:
% \section{Extension of the Framework}
% \begin{frame}{frame title}
%     frame contents
% \end{frame}




% % TODO:
% \subsection{Revisiting the Proof}



% \begin{frame}{Applied to the 1D Model}
%
%     \begin{itemize}
%         \item From the Proposition, we see that 1-var will decay
%         \item The behavior is then of course easily carried to the higher
%             dimensional models
%         \item But we of course care about the transient state
%     \end{itemize}
%
% \end{frame}

\section{ Over Now }


\begin{frame}[t, allowdisplaybreaks=4]{Bibliography}
    %, allowframebreaks=0.25]


    % \vspace{2em}
    \printbibliography


\end{frame}

\begin{frame}[t]{End}

    % \vspace{2em}
    \begin{itemize}
        \item Thank
        \item Questions?
    \end{itemize}

\end{frame}





\appendix
\section{Auxiliary}

\frame[noframenumbering]{\sectionpage}

\begin{frame}{Waddington's Figures}

    \begin{columns}%[T]

        \column{0.5\framewidth}
    \begin{figure}[htbp]
    \centering
    \includegraphics[width=0.9\textwidth]{epigenetic.png}
    \caption{
        Part of an Epigenetic Landscape.
    }
    \end{figure}

        \column{0.5\framewidth}
    \begin{figure}[htbp]
    \centering
    \includegraphics[width=1\textwidth]{factorsofepigenetic.png}
    \caption{
        The complex system of interactions underlying the epigenetic landscape.
    }
    \end{figure}

    \end{columns}

\end{frame}

\begin{frame}[t]{ Weird Outcomes }

    % \vspace{2em}
    \begin{itemize}
        \item Recalling the unexpected behavior in Figures 6.4 and 6.5 ( on page 36 )
    \end{itemize}

    \begin{tikzpicture}[overlay, remember picture]

        \node at (frame cs:-0.5, 0) {
                \includegraphics[width=0.5\textwidth,page=1]{ssa:ode-5-2:I19:C2fp:T17475890580172482.pdf}
            };

        \node at (frame cs: 0.5, 0) {
                \includegraphics[width=0.5\textwidth,page=2]{ssa:ode-5-2:I19:C2fp:T17475890580172482.pdf}
            };

    \end{tikzpicture}


\end{frame}

\subsection{My Model in 1-Dimension}
\begin{frame}{Simplest Novel Model}

    \begin{itemize}
        \item The one dimensional, i.e., just the GCN
        \item So to rectify that
    \end{itemize}

    \begin{columns}%[C]

        \column{0.5\framewidth}

        \begin{reactiontable}
            \reactionterm[\(k_{-1} x^2\)]{\(k_1 n x\)}{\(X+N\)}{\(2 X\)}{ 1 }
            \reactionterm{\(k_3 b x\)}{\(X+B\)}{\(B+N\)}{ -1 }
        \end{reactiontable}

        \column{0.4\framewidth}

        \begin{align*}
            A\of x &= k_{-1} x^2 + \pa{ k_3 b - k_1 n } x
            \\
            B\of x &= {1 \over \Omega} \pa{k_{-1} x^2 + \pa{ k_1 n + k_3 b } x}
        \end{align*}

    \end{columns}

\end{frame}

\begin{frame}{ Analytical Distribution }

    \begin{columns}

        \column{0.60\framewidth}

        \begin{itemize}
            \item Has an exponential distribution
                \begin{itemize}
                    \item Which is what I would hope for
                \end{itemize}
            \item Very dependent to changes in ratio of \(N\) to \(B\)
                \begin{itemize}
                    \item \(N \approx  B \implies\) distribution centered on \(x=0\)
                    \item \(N \not\approx B \implies\) distribution centered on \( (k_1n - k_3b) / k_{-1}\)
                \end{itemize}
            \item Would have made the rest of my life easier to see the
                problems from this version
        \end{itemize}

        \column{0.6\framewidth}

    \end{columns}

    \begin{tikzpicture}[overlay, remember picture]
        \node at (frame cs:0.55, 0.0) {
                \includegraphics[width=0.5\textwidth]
                {Novel1VarData.pdf}
            };
    \end{tikzpicture}

\end{frame}

\begin{frame}{The ODE System}

    \tikz[overlay] \node<1-> at (frame cs:-0.65,-1.35) {
        \begin{tikzpicture}[ scale = 1,
            draw = black!80,
            thick,
            shorten >=2pt,
            overlay
            ]

            \node[draw=black, circle] (c1) at (-1, 0) {\(c_1\)};
            \node[draw=black, circle] (c2) at ( 1, 0) {\(c_2\)};

            \node[draw=black] (m) at ( -1, -1.20) {\(m\)};
            \node[] (source) at ( -2, -1.20) {\(m_0\)};

            \draw[->] (c2) [bend left]to node[midway, below] {\(w_2\)} (c1);
            \draw[->] (c1) [bend left]to node[midway, above] {\(w_1\)} (c2);

            \draw[->] (c1)  [loop left]to node[midway, left]  {\(k_1\)} (c1);
            \draw[->] (c2) [loop right]to node[midway, right] {\(k_2\)} (c2);

            \draw[-|] (m) to node[pos=0.35] (mm) {} (c1);
            \draw[-|, rounded corners] (c2) |- node[midway,anchor=south west] {\(q_m\)} (m);
            % \draw[-|] (c2) to[out=270, in=0, looseness=2] node[midway,anchor=north west] {\(q_m\)} (m);

            \node[anchor=east] at (mm) {\(q_1\)};

            \draw[->] (source) -- (m);

    \end{tikzpicture}};


    \begin{columns}
        \column{0.35\framewidth}
        \vspace{8em}\par


        \column{0.5\framewidth}
        \begin{align*}
            \ode{c_1} &= k_1 c_1\left({n_1 - c_t \over n_1}\right) - w_1 c_1 + w_2 c_2 - q_1 m c_1 \nonumber \\
            \ode{c_2} &= k_2 c_2\left({n_2 - c_t \over n_2}\right) - w_2 c_2 + w_1 c_1\nonumber\\
            \ode{m}   &= m_0 - q_m m c_2 \nonumber
        \end{align*}

    \end{columns}

    \begin{itemize}
        \item Shown to be Lyapunov stable at the sole fixed point.
    \end{itemize}

\end{frame}

\begin{frame}{The ODE Dynamics}


    \begin{onlyenv}<1>
        \begin{figure}[htbp]
            \centering
            \begin{subfigure}{0.5\framewidth}
                \centering
                \includegraphics[width=1\textwidth,page=1]{ode:jesper:I15:C1:T1748197107286477.pdf}
                \caption{The ODE system with \(m_0 = 0.20\)}
            \end{subfigure}~
            \begin{subfigure}{0.5\framewidth}
                \centering
                \includegraphics[width=1\textwidth,page=1]{ode:jesper:I32:C1fp:T17481997939903586.pdf}
                \caption{The ODE system with \(m_0 = 10\)}
            \end{subfigure}
            \caption{ Two examples of the traditional ODE system with varying parameters }
        \end{figure}
    \end{onlyenv}

    \begin{onlyenv}<2>
        \begin{figure}[htbp]
            \centering
            \begin{subfigure}{0.5\framewidth}
                \centering
                \includegraphics[width=1\textwidth,page=1]{ssa:ODE8Par3Var:I34:C2:T17476639695953887.pdf}
            \end{subfigure}~
            \begin{subfigure}{0.5\framewidth}
                \centering
                \includegraphics[width=1\textwidth,page=2]{ssa:ODE8Par3Var:I34:C2:T17476639695953887.pdf}
            \end{subfigure}
            \caption{
                The stochastic version of the model simulated using
                Gillespie's SSA.
                The ODE system with \(m_0 = 10\)
            }
        \end{figure}
    \end{onlyenv}

    \begin{onlyenv}<3->
        \begin{figure}[htbp]
            \centering
            \begin{subfigure}{0.5\framewidth}
                \centering
                \includegraphics[width=1\textwidth,page=3]{ssa:ODE8Par3Var:I34:C2:T17476639695953887.pdf}
            \end{subfigure}~
            \begin{subfigure}{0.5\framewidth}
                \centering
                \includegraphics[width=1\textwidth,page=4]{ssa:ODE8Par3Var:I34:C2:T17476639695953887.pdf}
            \end{subfigure}
            \caption{
                The stochastic version of the model simulated using
                Gillespie's SSA.
                The ODE system with \(m_0 = 10\)
            }
        \end{figure}
    \end{onlyenv}

\end{frame}

\begin{frame}{ Stochastic Model Dynamics }

    \begin{onlyenv}<2->
        \begin{figure}[htbp]
            \centering
            \begin{subfigure}{0.5\framewidth}
                \centering
                \includegraphics[width=1\textwidth,page=1]{ssa:Novel5Par3Var:I26:C2:T17475033100262814.pdf}
            \end{subfigure}~
            \begin{subfigure}{0.5\framewidth}
                \centering
                \includegraphics[width=1\textwidth,page=1]{ssa:Novel5Par3Var:I22:C1:T17475031669379826.pdf}
            \end{subfigure}
            \caption{
                The stochastic version of the model simulated using
                Gillespie's SSA.
                The ODE system with \(m_0 = 10\)
            }
        \end{figure}
    \end{onlyenv}

\end{frame}

\begin{frame}{ Research Question Drafts }

    \vspace{2ex}
    \begin{columns}[T]

        \column{0.5\framewidth}

        \begin{enumerate}
            \item To what degree can modeled systems accurately determine the
                outcome of heteroresistence in bacteria
            \item Formalization of Waddington's epigenetic landscape through
                statistical manifolds
            \item How can one apply the Fokker-Planck equation to understand the
                evolution of distinct bacterial populations in relation to the
                bacterial property of heteroresistence of antibiotics
            \item How much insight can be obtained on heteroresistence through
                generalized models of ODE interactions
        \end{enumerate}

        \column{0.5\framewidth}

        \begin{enumerate}
            \item How can the FPE be used to model the stochastic dynamics of
                gene copy number in bacterial populations exhibiting heteroresistence.
            \item What are the key differences in prediction accuracy between
                ODE-based models and FPE-based models for GCN
                dynamics in bacteria
        \end{enumerate}

    \end{columns}



\end{frame}


\end{document}
